% ******************************************************************************
% ****************************** Custom Margin *********************************

% Add `custommargin' in the document class options to use this section
% Set {innerside margin / outerside margin / topmargin / bottom margin}  and
% other page dimensions
\ifsetCustomMargin
  \RequirePackage[left=37mm,right=30mm,top=35mm,bottom=30mm]{geometry}
  \setFancyHdr % To apply fancy header after geometry package is loaded
\fi

% *****************************************************************************
% ******************* Fonts (like different typewriter fonts etc.)*************

% Add `customfont' in the document class option to use this section

\ifsetCustomFont
  % Set your custom font here and use `customfont' in options. Leave empty to
  % load computer modern font (default LaTeX font).
  \RequirePackage{helvet}
\fi

% *****************************************************************************
% **************************** Custom Packages ********************************

% ************************* Algorithms and Pseudocode **************************

%\usepackage{algpseudocode}


% ********************Captions and Hyperreferencing / URL **********************

% Captions: This makes captions of figures use a boldfaced small font.
%\RequirePackage[small,bf]{caption}

\RequirePackage[labelsep=space,tableposition=top]{caption}
\renewcommand{\figurename}{Fig.} %to support older versions of captions.sty


% *************************** Graphics and figures *****************************

\usepackage{rotating}
%\usepackage{wrapfig}

% Uncomment the following two lines to force Latex to place the figure.
% Use [H] when including graphics. Note 'H' instead of 'h'
%\usepackage{float}
%\restylefloat{figure}

% Subcaption package is also available in the sty folder you can use that by
% uncommenting the following line
% This is for people stuck with older versions of texlive
%\usepackage{sty/caption/subcaption}
% \usepackage{subcaption}

% ********************************** Tables ************************************
\usepackage{booktabs} % For professional looking tables
\usepackage{colortbl}
\usepackage{multirow}
\usepackage{hhline}
\usepackage{longtable}
\usepackage{threeparttable}
\usepackage{subfigure}
\usepackage{adjustbox}

% \usepackage{minitoc}

% ***************************** Math and SI Units ******************************
\usepackage{framed,enumitem} 

\usepackage{siunitx} % use this package module for SI units
\usepackage[hyphens]{xurl} % better URL breaking at more positions
\usepackage{hyperref}
\hypersetup{
  colorlinks,
  citecolor=blue,
  linkcolor=magenta,
  urlcolor=cyan}

% Bibliography links (plainnat-fl writes \doi{...} for the doi field).
% (1) If doi is set: show one clickable DOI link — use full URL as-is when it is
%     already http(s); otherwise https://doi.org/{id} after stripping optional
%     doi:/DOI: prefix. (2) If doi is empty, plainnat-fl prints url only (format.url).
% Defined before the .bbl so \providecommand{\doi} in .bbl is skipped.
\makeatletter
\newcommand{\doi}[1]{%
  \begingroup
  \urlstyle{rm}%
  \edef\DOI@full{\detokenize{#1}}%
  \expandafter\DOI@emit\DOI@full\DOI@tail
  \endgroup
}
\def\DOI@tail{\DOI@tail}
% Split off first token: http(s) URLs keep full string #1#2; else resolve bare DOI.
% Use \pdfstrcmp (pdfTeX): \ifnum`#1=`h breaks when #1 is empty or not a plain char token.
\def\DOI@emit#1#2\DOI@tail{%
  \def\DOI@first{#1}%
  \ifnum\pdfstrcmp{\DOI@first}{h}=0
    \href{#1#2}{\url{#1#2}}%
  \else
    \ifnum\pdfstrcmp{\DOI@first}{H}=0
      \href{#1#2}{\url{#1#2}}%
    \else
      \DOI@bare{#1#2}%
    \fi
  \fi
}
% Strip leading doi: / DOI: (optional); then always use resolver URL for display.
\def\DOI@bare#1{%
  \DOI@A #1\DOI@@@%
}
\def\DOI@A doi:#1\DOI@@@{\DOI@link{#1}}
\def\DOI@A DOI:#1\DOI@@@{\DOI@link{#1}}
\def\DOI@A #1\DOI@@@{\DOI@link{#1}}
\def\DOI@link#1{%
  \href{https://doi.org/#1}{\url{https://doi.org/#1}}%
}
\makeatother

% \usepackage[square,numbers]{natbib}
\usepackage{graphicx}
\usepackage{amsfonts,amssymb,amsmath,amsthm, mathtools}
\usepackage[version=4]{mhchem}
\DeclareGraphicsExtensions{.pdf,.png,.jpg,.jpeg}
\usepackage{multicol}
\usepackage{subfigure}
\setlength{\arrayrulewidth}{0.3mm}
\usepackage{geometry}
 \geometry{
 a4paper,
 % total={170mm,257mm},
 % left=20mm,
 top=25mm,
}
% Uniform paragraph spacing: twoside + \flushbottom (book default) stretches
% vertical glue on short pages; \raggedbottom lets pages end short instead.
% Without that stretch, book class keeps \parskip=0pt, so paragraphs look tight;
% set a fixed parskip (no plus/minus) so gaps stay even and do not vary by page.
\raggedbottom
\setlength{\parskip}{0.45\baselineskip}
% Optional: use block paragraphs (no first-line indent) instead of the default
% indented style: uncomment the next line.
% \setlength{\parindent}{0pt}
% A blank line after a display starts a new paragraph. Keep that paragraph
% flush: the formula does not begin a new indented block.
\usepackage{noindentafter}
\NoIndentAfterEnv{equation}
\NoIndentAfterEnv{equation*}
\NoIndentAfterEnv{align}
\NoIndentAfterEnv{align*}
\NoIndentAfterEnv{gather}
\NoIndentAfterEnv{gather*}
\NoIndentAfterEnv{multline}
\NoIndentAfterEnv{multline*}
\usepackage{algorithm} 
\usepackage[algo2e]{algorithm2e} 
\usepackage{algpseudocode}
\SetKw{KwBy}{by}
% \usepackage{algorithm}
% \usepackage{algpseudocode}

%Options: Sonny, Lenny, Glenn, Conny, Rejne, Bjarne, Bjornstrup
\usepackage[Bjornstrup]{fncychap}

% ******************************************************************************
% % Format pages for PART
% Load xcolor to replace the old color package from the class file
% xcolor automatically replaces color when loaded, but we ensure compatibility
% \usepackage{xcolor}
\usepackage[dvipsnames]{xcolor}
\definecolor{RED}{rgb}{1,0,0} % Define RED as an alias for red in case something uses uppercase
% Load tocloft before etoc to avoid conflicts
% Prevent counter redefinition errors by patching \newcounter
\makeatletter
% Save original \newcounter
\let\@@newcounter\newcounter
% Redefine \newcounter to check if counter exists first
\renewcommand{\newcounter}[1]{%
  \@ifundefined{c@#1}{%
    \@@newcounter{#1}%
  }{%
    % Counter already exists, do nothing to avoid error
  }%
}
\usepackage{tocloft}
% Restore original \newcounter after tocloft loads
\let\newcounter\@@newcounter
\makeatother
\usepackage{etoc}
% Packages for PART and EPIGRAPH
\usepackage{epigraph}
% indentafter: indent the first paragraph after a chapter or section title,
% same as every other paragraph.
\usepackage[indentafter]{titlesec}
\titleformat{\part}[display]
  {\filleft\fontsize{40}{40}\selectfont\scshape}
  {\fontsize{90}{90}\selectfont\thepart}
  {20pt}
  {\thispagestyle{epigraph}}

\setlength\epigraphwidth{.6\textwidth}
%%%%


% ******************************************************************************
% created new environemnt to add Catalan Abstract
\newenvironment{catalanabstract} {
  \cleardoublepage
  \setsinglecolumn
  \chapter*{\centering \Large Resum}
  \thispagestyle{empty}
}
% ******************************************************************************
% ******************************************************************************
% created new environemnt to add Spanish Abstract

\newenvironment{spanishabstract} {
  \cleardoublepage
  \setsinglecolumn
  \chapter*{\centering \Large Resumen}
  \thispagestyle{empty}
}


% \makeglossaries


% ******************************* Line Spacing *********************************

% Choose linespacing as appropriate. Default is one-half line spacing as per the
% University guidelines

% \doublespacing
% \onehalfspacing
% \singlespacing

% Borrar si és necessari. Afegit per reduir \hbox overfull warnings.
% Allow more flexible line breaking to reduce overfull hbox warnings
\setlength{\emergencystretch}{3em}


% ************************ Formatting / Footnote *******************************

% Don't break enumeration (etc.) across pages in an ugly manner (default 10000)
%\clubpenalty=500
%\widowpenalty=500

%\usepackage[perpage]{footmisc} %Range of footnote options


% *****************************************************************************
% *************************** Bibliography  and References ********************

%\usepackage{cleveref} %Referencing without need to explicitly state fig /table

% Add `custombib' in the document class option to use this section
\ifuseCustomBib
   \RequirePackage[square, sort, numbers, authoryear]{natbib} % CustomBib

% If you would like to use biblatex for your reference management, as opposed to the default `natbibpackage` pass the option `custombib` in the document class. Comment out the previous line to make sure you don't load the natbib package. Uncomment the following lines and specify the location of references.bib file

%\RequirePackage[backend=biber, style=numeric-comp, citestyle=numeric, sorting=nty, natbib=true]{biblatex}
%\bibliography{References/references} %Location of references.bib only for biblatex

\fi

% changes the default name `Bibliography` -> `References'
\renewcommand{\bibname}{References}


% *****************************************************************************
% *************** Changing the Visual Style of Chapter Headings ***************
% This section on visual style is from https://github.com/cambridge/thesis

% Uncomment the section below. Requires titlesec package.

%\RequirePackage{titlesec}
%\newcommand{\PreContentTitleFormat}{\titleformat{\chapter}[display]{\scshape\Large}
%{\Large\filleft{\chaptertitlename} \Huge\thechapter}
%{1ex}{}
%[\vspace{1ex}\titlerule]}
%\newcommand{\ContentTitleFormat}{\titleformat{\chapter}[display]{\scshape\huge}
%{\Large\filleft{\chaptertitlename} \Huge\thechapter}{1ex}
%{\titlerule\vspace{1ex}\filright}
%[\vspace{1ex}\titlerule]}
%\newcommand{\PostContentTitleFormat}{\PreContentTitleFormat}
%\PreContentTitleFormat


% ******************************************************************************
% ************************* User Defined Commands ******************************
% ******************************************************************************

% *********** To change the name of Table of Contents / LOF and LOT ************

%\renewcommand{\contentsname}{My Table of Contents}
%\renewcommand{\listfigurename}{My List of Figures}
%\renewcommand{\listtablename}{My List of Tables}

% \renewcommand{\listfigurename}{List of Figures}
% \renewcommand{\listtablename}{List of Tables}
% \renewcommand{\contentsname}{Contents}
% \renewcommand{\listalgorithmname}{List of Algorithms}


% ********************** TOC depth and numbering depth *************************

\setcounter{secnumdepth}{3}
\setcounter{tocdepth}{3}

% ********************** List of Figures and Tables spacing ********************
% Increase spacing between number and title in list of figures and tables
% (tocloft is loaded earlier, before etoc)
% Add space after the number to prevent superposition with the title
\renewcommand{\cftfigaftersnum}{\quad}  % Add space after figure number (quad = 1em)
\renewcommand{\cfttabaftersnum}{\quad}  % Add space after table number (quad = 1em)
\setlength{\cftfignumwidth}{3.6em}  % Width for figure numbers
\setlength{\cfttabnumwidth}{3.6em}  % Width for table numbers


% ******************************* Nomenclature *********************************

% To change the name of the Nomenclature section, uncomment the following line

%\renewcommand{\nomname}{Symbols}


% ********************************* Appendix ***********************************

% The default value of both \appendixtocname and \appendixpagename is `Appendices'. These names can all be changed via:

%\renewcommand{\appendixtocname}{List of appendices}
%\renewcommand{\appendixname}{Appndx}

% ******************************** Draft Mode **********************************

% Uncomment to disable figures in `draftmode'
%\setkeys{Gin}{draft=true}  % set draft to false to enable figures in `draft'

% These options are active only during the draft mode
% Default text is "Draft"
%\SetDraftText{DRAFT}

% Default Watermark location is top. Location (top/bottom)
%\SetDraftWMPosition{bottom}

% Draft Version - default is v1.0
%\SetDraftVersion{v1.1}

% Draft Text grayscale value (should be between 0-black and 1-white)
% Default value is 0.75
%\SetDraftGrayScale{0.8}


%% Todo notes functionality
%% Uncomment the following lines to have todonotes.

%\ifsetDraft
%	\usepackage[colorinlistoftodos]{todonotes}
%	\newcommand{\mynote}[1]{\todo[author=kks32,size=\small,inline,color=green!40]{#1}}
%\else
%	\newcommand{\mynote}[1]{}
%	\newcommand{\listoftodos}{}
%\fi

% Example todo: \mynote{Hey! I have a note}




